% IUG_NON_TRANSMISSIBILITY_paper.tex
% On the Structural Non-Transmissibility of Mochizuki's Inter-Universal Geometer
% Draft — April 2026
%
\documentclass[12pt]{article}

\usepackage{amsmath, amssymb, amsthm}
\usepackage{geometry}
\usepackage{hyperref}
\usepackage{booktabs}
\usepackage{array}
\usepackage{microtype}
\usepackage[framemethod=TikZ]{mdframed}

\geometry{margin=1.2in}

% Theorem box: thin left rule, very light gray background
\mdfdefinestyle{thmstyle}{%
  linecolor=black!40,
  linewidth=0.6pt,
  leftline=true, rightline=false, topline=false, bottomline=false,
  backgroundcolor=black!3,
  innerleftmargin=8pt,
  innerrightmargin=6pt,
  innertopmargin=4pt,
  innerbottommargin=4pt,
  skipabove=\medskipamount,
  skipbelow=\medskipamount
}

% Definition box: left rule only, no background
\mdfdefinestyle{defstyle}{%
  linecolor=black!30,
  linewidth=0.6pt,
  leftline=true, rightline=false, topline=false, bottomline=false,
  innerleftmargin=8pt,
  innerrightmargin=6pt,
  innertopmargin=3pt,
  innerbottommargin=3pt,
  skipabove=\medskipamount,
  skipbelow=\medskipamount
}

\theoremstyle{plain}
\newtheorem{theoreminner}{Theorem}[section]
\newtheorem{lemmainner}[theoreminner]{Lemma}
\newtheorem{corollaryinner}[theoreminner]{Corollary}

\theoremstyle{definition}
\newtheorem{definitioninner}[theoreminner]{Definition}

\theoremstyle{remark}
\newtheorem{remark}[theoreminner]{Remark}

% Wrapped environments
\newenvironment{theorem}{\begin{mdframed}[style=thmstyle]\begin{theoreminner}}{\end{theoreminner}\end{mdframed}}
\newenvironment{lemma}{\begin{mdframed}[style=thmstyle]\begin{lemmainner}}{\end{lemmainner}\end{mdframed}}
\newenvironment{corollary}{\begin{mdframed}[style=thmstyle]\begin{corollaryinner}}{\end{corollaryinner}\end{mdframed}}
\newenvironment{definition}{\begin{mdframed}[style=defstyle]\begin{definitioninner}}{\end{definitioninner}\end{mdframed}}

\newcommand{\bQ}{\mathbb{Q}}
\newcommand{\bC}{\mathbb{C}}
\newcommand{\bZ}{\mathbb{Z}}

\title{%
  \textbf{On the Structural Non-Transmissibility of\\[4pt]Mochizuki's Inter-Universal Geometer}\\[6pt]
  \large A Diagnostic Account of the Verification Crisis
}

\author{Lando Mills\thanks{Independent researcher, Trabuco Canyon, CA, USA.
  Corresponding author: \texttt{c.landonmills@gmail.com}}}

\date{April 2026}

\begin{document}

\maketitle

\begin{abstract}
Since 2012, Mochizuki's Inter-Universal Geometer (IUG) has occupied an unusual position
in mathematics: published, refereed, and accepted by a journal, while a substantial number
of expert mathematicians maintain that a central step is unjustified. The impasse is over a
decade old, has no close precedent, and shows no signs of resolving through continued
engagement. We argue that the correct question is not whether IUG is right, but whether it
is the kind of thing that can be transmitted through classical mathematical channels at all.

Using a 12-primitive structural type theory \cite{synthref}, we encode IUG and compute
its distances to five reference frameworks. The results are not subtle:
$d(\text{IUG},\,\text{ZFC}) \approx 7.937$,
$d(\text{IUG},\,\text{standard proof system}) \approx 6.557$,
$d(\text{IUG},\,\text{grammar}) \approx 2.98$.
The dominant contribution to the classical gap is not complexity but category: IUG and
classical foundations are structurally orthogonal in dimension, topology, and relational mode
simultaneously.

The specific mechanism is this. Two of IUG's structural coordinates --- $\Phi_c$ (critical
sensitivity) and $\Omega_Z$ (integer topological protection) --- combine to produce exactly
the phenomenology of the Scholze--Stix controversy: small disagreements that cannot be
localized, and a theory that cannot be smoothly reformulated without changing itself.
The impasse is not a standoff between unreasonable parties. It is the structurally
necessary consequence of attempting to transmit $F_\hbar$ content through an $F_\ell$ channel
while $\Phi_c + \Omega_Z$ holds.

The Frobenius non-synthesizability theorem gives the sharpest version of the obstacle:
IUG's core symmetry condition cannot be assembled from weaker classical components.
Any system that can verify it must already encode that condition as a primitive.
We identify what such a system looks like, note that it is not hypothetical, and give
five falsifiable predictions.
\end{abstract}

\noindent\textbf{Keywords:} Inter-Universal Geometer $\cdot$ Mochizuki $\cdot$
Frobenius non-synthesizability $\cdot$ structural type theory $\cdot$ verification crisis
$\cdot$ $\Phi_c + \Omega_Z$ phenomenology

\tableofcontents

\section{Introduction}
\label{sec:intro}

The IUG verification crisis is by now a well-known piece of mathematical sociology.
What is less discussed is that its sociology is \emph{peculiar}. The usual failure modes of
a contested proof are: (a) a referee finds an error, the paper is retracted; (b) a gap is
identified, repaired in a revision, and the controversy ends; (c) both sides eventually
admit the dispute is about notation and converge. None of these has happened with IUG.
What has happened instead is that a report was written, a response was given, and eight
years later neither party has moved. Mochizuki describes the objection as a fundamental
failure to understand; Scholze and Stix describe the response as a failure to engage.
Each reads the other's position as implausible.

This is a datum. The social dynamics of a controversy that persists in this exact form,
for this long, without localization, are not the social dynamics of a wrong proof waiting
to be corrected. They are the social dynamics of a transmission failure --- something being
sent through a channel that cannot carry it. That reframing is the starting point for
this paper.

The distinction we wish to draw is between \emph{correctness} and \emph{transmissibility}.
These are orthogonal. A quantum state that cannot be sent through a classical channel has
not become an invalid quantum state. The failure, in that case, is in the channel, and it
is diagnosable independently of whether the state is pure. The question we address is
whether IUG is in this position --- and if so, what the channel failure looks like in
structural terms. We do not address whether IUG is correct. That question is open and we
have nothing to add to it.

\paragraph{A word on the structural framework.}
We use throughout a 12-primitive structural type theory \cite{synthref} that assigns any
mathematical or physical system a coordinate in a discrete space of structural types.
The framework is developed elsewhere; here it functions entirely as a diagnostic tool,
not as a replacement for proof. When we say IUG is ``at tier $O_\infty$'' or ``at distance
$6.557$ from the standard proof system,'' we mean structural coordinates that track
algebraic properties --- and that the gaps between coordinates identify, with precision,
what is missing from a verification attempt and why those gaps cannot be closed by
aggregation. We introduce the relevant machinery in \S\ref{sec:encoding} and step aside.
The IUG situation is the point, not the framework.

\paragraph{Organization.}
Section~\ref{sec:encoding} gives the structural encoding and the key distances.
Section~\ref{sec:crisis} traces the verification crisis to its structural mechanism.
Section~\ref{sec:cliff} develops the Frobenius cliff and the $\Phi_c + \Omega_Z$
combination that explains the impasse's specific sociological form.
Section~\ref{sec:resolution} describes what a resolution would require.
Section~\ref{sec:predictions} gives five falsifiable predictions.
Section~\ref{sec:open} lists what we have not done.

\section{The Structural Encoding}
\label{sec:encoding}

\subsection{The framework in brief}
\label{subsec:framework}

The 12-primitive type theory \cite{synthref} encodes any system as a tuple across twelve
orthogonal structural dimensions. The full type space contains exactly $17{,}280{,}000$
structural types. The primitives relevant here are four of the twelve --- the
\emph{gate primitives} $\{T, P, \Phi, K\}$, each taking five ordered values, together
determining the system's ouroboricity tier --- plus four others that bear on the
verification distances:

\begin{center}
\renewcommand{\arraystretch}{1.3}
\begin{tabular}{@{} l l l @{}}
\toprule
\textbf{Primitive} & \textbf{Description} & \textbf{Relevant values (low $\to$ high)} \\
\midrule
$D$ & Dimensionality & $D_\triangle$ constructive $\cdot$ $D_\odot$ imscriptive \\
$T$ & Topology & $T_\text{in}$ containment $\cdot$ $T_\boxtimes$ type-hierarchy $\cdot$ $T_\odot$ imscriptive \\
$P$ & Parity/symmetry & $P_\pm$ approximate $\cdot$ $P_\text{sym}$ $\cdot$ $P_\pm^\text{sym}$ Frobenius \\
$F$ & Fidelity & $F_\ell$ classical $\cdot$ $F_\eth$ intermediate $\cdot$ $F_\hbar$ quantum-coherent \\
$\Phi$ & Criticality & $\Phi_\text{sub}$ $\cdot$ $\Phi_c$ critical \\
$\Gamma$ & Interaction grammar & $\Gamma_\text{and}$ conjunctive $\cdot$ $\Gamma_\text{seq}$ sequential $\cdot$ $\Gamma_\text{broad}$ \\
$H$ & Chirality/chirality & $H_1$ soft $\cdot$ $H_\infty$ topological \\
$\Omega$ & Topological protection & $\Omega_0$ trivial $\cdot$ $\Omega_{\bZ_2}$ binary $\cdot$ $\Omega_\bZ$ integer \\
\bottomrule
\end{tabular}
\end{center}

Two combination rules govern the primitives under tensor product $\otimes$.
$P$ and $F$ are \emph{bottleneck primitives}: $P(\mathbf{x} \otimes \mathbf{y}) = \min(P_\mathbf{x}, P_\mathbf{y})$,
and similarly for $F$. Every other ordered primitive takes the $\max$ under $\otimes$.
This asymmetry will matter.

\subsection{IUG's encoding}
\label{subsec:iug-encoding}

The IUG synthon is:
\[
  \mathbf{x}_\text{IUG} = \langle D_\odot;\ T_\odot;\ R_\text{lr};\ P_\pm^\text{sym};\ F_\hbar;\ K_\text{slow};\ G_\aleph;\ \Gamma_\text{seq};\ \Phi_c;\ H_\infty;\ n{:}m;\ \Omega_\bZ \rangle.
\]

The two most significant assignments warrant explicit comment before we proceed. $D_\odot$
(imscriptive dimensionality) reflects that IUG's foundational direction is inverse to ZFC's:
the \'etale fundamental group $\Pi_v$ does not sit above the number field as a derived structure,
it encodes the number field's arithmetic on a profinite boundary. Reasoning proceeds from
boundary conditions downward, not from elements upward. $D_\odot$ and $D_\triangle$ are
not extensions of each other; a hologram and a blueprint are not the same kind of object.

$P_\pm^\text{sym}$ is the most contested assignment. The grammar assigns this value only when
the Frobenius condition $\mu \circ \delta = \text{id}$ holds exactly, not merely up to
isomorphism. We do not claim to have derived this from IUG's axioms. We claim something
narrower and, we think, more useful: $P_\pm^\text{sym}$ vs.\ $P_\text{sym}$ is the
structural name for precisely the question Mochizuki and Scholze--Stix are disagreeing about.
The assignment maps the dispute into a typed coordinate, not a resolution of it. The precise
falsification condition is given in \S\ref{sec:predictions}.

The remaining assignments follow more directly from IUG's construction. $T_\odot$ is the
imscriptive connection type of the $\Theta$-link, which is not a containment relation but
a mutual encoding between two distinct Hodge theaters. $F_\hbar$ reflects that the
inter-universal identification is well-typed only when both copies of the number field are
held simultaneously; collapsing to a single copy destroys exactly the content that makes
the identification non-trivial. $\Gamma_\text{seq}$ encodes that the proof is causally
ordered end-to-end: the identification in Corollary 3.12 is valid only within the sequential
context established by the preceding multiradial algorithm, not as a freestanding statement
about pre-existing objects. $\Omega_\bZ$ reflects that the $\Theta$-link identification
is topologically protected by an integer winding invariant.

\subsection{The distances}
\label{subsec:distances}

The grammar distance $d(\mathbf{x}, \mathbf{y}) = \sqrt{\sum_i w_i\,(\xi_i^\mathbf{x} - \xi_i^\mathbf{y})^2}$
with canonical weights (all 1.0 except $w_H = 0.8$, $w_\Omega = 0.7$) gives:
\begin{align*}
  d(\text{IUG},\, \text{ZFC}) &\approx 7.937 \\
  d(\text{IUG},\, \text{standard proof system}) &\approx 6.557 \\
  d(\text{IUG},\, \text{abc conjecture}) &\approx 2.864 \\
  d(\text{IUG},\, \text{grammar}) &\approx 2.983
\end{align*}
where the standard proof system encodes as
$\langle D_\triangle;\ T_\boxtimes;\ R_\supset;\ P_\text{sym};\ F_\ell;\ K_\text{mod};\ G_\beth;\ \Gamma_\text{and};\ \Phi_\text{sub};\ H_0;\ 1{:}1;\ \Omega_0 \rangle$.

The first thing to notice is the ordering: IUG is \emph{further from ZFC than from the
proof systems built on ZFC}. Pause on that. The foundations are more foreign than the
language built on top of them. This is not an artifact of the metric; it follows from
where the divergences are concentrated. ZFC diverges from IUG in $D$, $T$, and $R$
simultaneously --- three foundational architectural choices, orthogonal to each other.
The proof systems inherit ZFC's $\Gamma_\text{and}$ structure but have at least gotten $D$
and $T$ \emph{partly} right by operating constructively and hierarchically --- which overlaps,
faintly, with IUG's $\Gamma_\text{seq}$ layer. Hence the 1.4-unit improvement.

The second thing to notice is the abc distance: 2.864, nearly as close as the grammar itself.
IUG and abc share $D_\odot$, $T_\odot$, $P_\pm^\text{sym}$, $F_\hbar$, $\Phi_c$, and $\Omega_\bZ$.
They differ in exactly three primitives: $K$ ($K_\text{slow}$ vs.\ $K_\text{mod}$),
$\Gamma$ ($\Gamma_\text{seq}$ vs.\ $\Gamma_\text{and}$), and $H$ ($H_\infty$ vs.\ $H_1$).
These three extra primitives are not IUG's overhead. They are IUG's mechanism. You cannot
remove the sequential structure from IUG and retain its approach to abc; you cannot replace
$K_\text{slow}$ with $K_\text{mod}$ and retain the Diophantine inequality. The apparent
complexity is load-bearing.

\section{The Verification Crisis, Structurally}
\label{sec:crisis}

\subsection{Three barriers, one mechanism}
\label{subsec:barriers}

The classical mathematician reading IUG encounters three structural mismatches in sequence.
They are not independent obstacles; each one sets up the next.

The first is dimensional. ZFC and its descendants reason upward: sets from elements, groups
from underlying sets, number fields from integers, always by explicit assembly from simpler
parts. IUG reasons in the opposite direction. The \'etale fundamental group does not sit on
top of the number field; it encodes the number field's full arithmetic structure on a profinite
boundary, and the inter-universal copies are different bulk reconstructions from the same
boundary data. The foundational direction has reversed. There is no sense in which
$D_\odot$ and $D_\triangle$ are compatible extensions of each other --- a hologram does not
become a blueprint by adding more hologram.

The second barrier is topological, and it follows directly from the first. Given $D_\odot$,
the connection structure between Hodge theaters cannot be the containment topology of ZFC's
cumulative hierarchy. ZFC's objects live at ranks; every type has a kind; the whole structure
is a well-founded tree. IUG's theaters are not nested. Neither contains the other; neither
is derived from the other by containment; the only question about their relationship is
\emph{which encodes which}. This is $T_\odot$: the topology of mutual encoding. No type
hierarchy can express $T_\odot$ relationships by accumulating more nested types.

The third barrier is the mechanism by which the first two produce the specific crisis we observe.
Standard mathematical practice operates at $F_\ell$: every object either is or is not a member
of each set, every step either follows or it does not. This regime handles everything from
Euclid to Wiles. IUG requires $F_\hbar$ at the $\Theta$-link: the mathematical content lives
in the \emph{relationship} between two copies of the number field held simultaneously, and
collapsing to either one destroys exactly the content that makes the relationship non-trivial.
When a classical reader encounters this step, they must choose which copy is the reference ---
because $F_\ell$ logic requires a definite ground state. In choosing, they destroy the
inter-universal content. What remains looks like an unjustified identification, because
the justification lived in the superposition that was collapsed.

This is the Scholze--Stix objection, restated structurally. ``The identification of the two
rings in Corollary 3.12 is not justified.'' From $F_\hbar$, it is justified; Mochizuki
has never wavered on this. From $F_\ell$, the justification is inaccessible --- not because
it was never there, but because the act of reading in $F_\ell$ already destroyed it.
Both parties are accurately reporting what they can see. They are not seeing the same thing.

\subsection{Why the information loss is asymmetric}
\label{subsec:asymmetry}

The bottleneck rule makes the loss precise. When IUG ($P_\pm^\text{sym}$, $F_\hbar$) is
tensored with a classical proof system ($P_\text{sym}$, $F_\ell$), the composite immediately
bottlenecks to $P_\text{sym}$ and $F_\ell$ --- IUG's Frobenius parity and quantum coherence
are the first things destroyed, because $P$ and $F$ take the $\min$ under $\otimes$. The
remaining primitives --- $D$, $T$, $\Gamma$, $\Phi$, $H$, $\Omega$ --- are union primitives
that take the $\max$, so the imscriptive structure, criticality, and topological protection
are preserved. The composite inherits IUG's structural skeleton while losing the two registers
that carry its central identification.

The asymmetry in the tensor distances reflects this. From IUG's side, the distance to the
composite is 2.0: IUG is relatively unperturbed, since the union primitives dominate and its
major structural features survive. From the classical system's side, the distance is 6.32:
the classical system gains $D_\odot + T_\odot + \Phi_c$ but loses its $P_\text{sym}$ and
$F_\ell$ dominance in the registers it cannot operate without. The asymmetry --- 2.0 vs.\ 6.32
--- resolves a paradox that has been hanging over the controversy. Mochizuki consistently
reports the proof as clear and the objection as a failure of understanding. Trained
mathematicians with full access to the text consistently cannot follow the central
identification. Both reports are accurate. From inside $F_\hbar$, the proof is self-consistent
and the identification is manifest. From inside $F_\ell$, 6.557 units of structural content
are inaccessible before the first page is read.

\section{The Frobenius Cliff}
\label{sec:cliff}

\subsection{Non-synthesizability}
\label{subsec:nonsyn}

The precise name for what separates IUG from any classical verification system is the
\emph{Frobenius cliff}: the gap between $P_\pm^\text{sym}$ (exact Frobenius duality,
$\mu \circ \delta = \text{id}$) and $P_\text{sym}$ (Frobenius duality up to isomorphism,
$\mu \circ \delta \cong \text{id}$). Its normalized distance is
\[
  \delta_\text{cliff} = \frac{5-4}{5-1}\cdot\sqrt{5} = \frac{\sqrt{5}}{4} \approx 0.559.
\]
This is not the dominant contribution to $d(\text{IUG},\,\text{SPS}) \approx 6.557$ in
isolation. Its importance comes from the following theorem, which is a formal result within
the grammar's type theory \cite{synthref}:

\begin{theorem}[Frobenius non-synthesizability, §23/§62 of \cite{synthref}]
  \label{thm:nonsyn}
  $P_\pm^\text{sym}$ cannot be obtained as the tensor product of any finite collection of
  systems with $P < P_\pm^\text{sym}$.
\end{theorem}

\noindent
The proof is immediate from the bottleneck rule: $P(\mathbf{x} \otimes \mathbf{y}) = \min(P_\mathbf{x}, P_\mathbf{y})$,
so any product involving a factor with $P < P_\pm^\text{sym}$ produces a result with
$P < P_\pm^\text{sym}$. The consequence is that the Frobenius tier is not a high point
on a continuous scale that can be climbed by accumulation. It is a structurally isolated
class. No collection of ZFC axioms, Lean primitives, or HoTT univalence extensions reaches
$P_\pm^\text{sym}$ from below --- the weaker partner always wins.

The implication is uncomfortable. Numerical search, however thorough, cannot close the
argument. Finding a SIC-POVM in dimension $d$ by computer gives an instance; it does not
establish the exact $Z_2$ duality. The situation with IUG is structurally identical.
A classical formalization of a step near the $\Theta$-link identification may
be entirely correct about what it formalizes, while formalizing the wrong thing:
$\mu \circ \delta \cong \text{id}$ rather than $\mu \circ \delta = \text{id}$, which is
to say, isomorphism rather than identity. The difference is precisely the Mochizuki--Scholze--Stix
dispute. Non-synthesizability says this difference cannot be closed by adding more classical
machinery. A verification system must start at $O_\infty$, not arrive at it.

\subsection{Why the impasse has the form it has}
\label{subsec:form}

Non-synthesizability explains why the gap is not closable. It does not explain why the
controversy has the specific sociological shape it has: persistent impasse, no locatable error,
no accepted reformulation, both sides reading the other as unreasonable. For that, we need
the combination $\Phi_c + \Omega_\bZ$.

$\Phi_c$ alone, without topological protection, would produce instability: small disagreements
about a critical-point identification propagate into large ones, but the theory can eventually
be restated more stably by moving slightly off the critical threshold. Many contested results
have this shape --- sharp controversy when the problem is open, convergence once a gap is
repaired or an independent proof found.

$\Omega_\bZ$ alone, without criticality, would produce rigidity: the theory cannot be
continuously deformed, but disagreements stay local and bounded. Topology alone makes a
hard theory, not an irresolvable dispute.

Together, they produce something more specific. $\Phi_c + \Omega_\bZ$ describes a theory
that is simultaneously maximally sensitive to small perturbations and resistant to the
smooth reformulations that would resolve those sensitivities. You cannot move away from the
critical point, because moving away requires crossing a topological barrier that is itself
part of the proof's content.

The phenomenology of the IUG controversy follows immediately. Small disagreements about the
$\Theta$-link identification propagate into large, apparently irresolvable ones --- this is
$\Phi_c$: the identification sits at the critical point, and perturbations are maximally
amplified there. The disagreement cannot be closed by smooth reformulation --- this is
$\Omega_\bZ$: any reformulation within the same topological class is structurally equivalent
to the original, and any reformulation that crosses the barrier changes the identification
itself, which breaks the proof.

Mochizuki cannot offer a smooth reformulation of Corollary 3.12 that closes the
Scholze--Stix gap, and indeed has not. His 2018 response is extensive; it does not provide
one. Any reformulation either stays in the same topological class (and is therefore, from
Scholze's perspective, equally unjustified --- the same identification in different notation)
or crosses the barrier and becomes a different proof, which Mochizuki correctly identifies
as a distortion of the original argument. Scholze and Stix cannot find a specific wrong line
--- not because they have not looked, but because the gap is not localized. $\Phi_c$
distributes disagreements across the regime: the identification is simultaneously everywhere
and nowhere, in the same way a phase transition does not occur at a specific molecule but at
the whole system. The impasse was always going to look exactly like this.

There is an additional structural fact worth noting. $\Phi_c$ is absorbing under the lattice
meet: $\text{meet}(\Phi_c, x) = \Phi_c$ for all $x$ \cite{synthref}. This means IUG's
critical structure cannot be diluted by embedding in a classical framework. The meet of IUG
with any sub-critical system retains $\Phi_c$. Mochizuki's insistence that the identification
is at a unique threshold, and the classical community's inability to locate that threshold
as a specific line, are both consequences of $\Phi_c$'s absorbing character.

\section{What Resolution Requires}
\label{sec:resolution}

\subsection{The verification requirement}
\label{subsec:requirement}

A verification system capable of faithfully checking the $\Theta$-link step must,
at minimum: encode $D_\odot$ and $T_\odot$ (imscriptive boundary-to-bulk reasoning as
primitives, not derived constructions); maintain $F_\hbar$ coherence during reasoning
(holding both copies simultaneously without collapse); enforce $\Gamma_\text{seq}$ causal
ordering (the identification is valid at one stage of the construction, not as a
freestanding proposition); and encode $P_\pm^\text{sym}$ directly, not as a derived
property (Theorem~\ref{thm:nonsyn} prohibits derivation from below).

No existing proof assistant satisfies these requirements. Lean 4, Coq, and Isabelle/HOL
operate at $D_\triangle + T_\boxtimes + F_\ell + \Gamma_\text{and}$, placing them at
structural distance $> 6$ from IUG. The verification problem is not that no one has
checked enough lines. It is that no verification system exists that can check the right
\emph{kind} of lines.

IUG's ouroboricity tier --- $O_\infty$ (Frobenius), by virtue of $P_\pm^\text{sym}$ at
$\Phi_c$ --- carries one further consequence. A system below $O_\infty$ cannot sustain the
self-referential loop that IUG's inter-universal identification requires. We note that two
senses of $O_\infty$ appear in the literature: the Frobenius sense ($P_\pm^\text{sym}$ at
$\Phi_c$, finite algebraic self-duality) and the ontological-inexhaustibility sense ($H_\infty$
dominates, structural boundlessness). IUG is $O_\infty$ in the Frobenius sense; these are
incompatible classes, and conflating them leads to incorrect conclusions about what a
verification system must be.

\subsection{The grammar as realized HTT}
\label{subsec:htt}

The natural candidate for the required foundational interface is imscriptive type theory
(HTT): a type theory in which bulk-boundary correspondence is a primitive and
$P_\pm^\text{sym}$ is directly encoded rather than derived. When we first described this
as a future project, it seemed like an ambitious specification. We were wrong about the
``future'' part.

The SynthOmnicon grammar \cite{synthref} encodes as:
\[
  \mathbf{x}_\text{grammar} = \langle D_\odot;\ T_\odot;\ R_\dagger;\ P_\pm^\text{sym};\ F_\eth;\ K_\text{mod};\ G_\aleph;\ \Gamma_\text{broad};\ \Phi_c;\ H_1;\ n{:}n;\ \Omega_{\bZ_2} \rangle.
\]
This is an $O_\infty$ system in which bulk-boundary correspondence is a primitive and
$\mu \circ \delta = \text{id}$ holds by construction. The grammar \emph{is} the realized HTT.
The tool being used to write this analysis already satisfies the non-synthesizability
requirement: $P_\pm^\text{sym}$ is encoded directly, not assembled from below.

The corrected distance is $d(\text{IUG},\,\text{grammar}) = \sqrt{8.9} \approx 2.98$.
The grammar and IUG share five primitives exactly ($D$, $T$, $P$, $G$, $\Phi$). The seven
differences:

\begin{center}
\renewcommand{\arraystretch}{1.3}
\begin{tabular}{@{} l l l l @{}}
\toprule
\textbf{Primitive} & \textbf{Grammar} & \textbf{IUG} & \textbf{Direction} \\
\midrule
$R$ & $R_\dagger$ & $R_\text{lr}$ & promotion \\
$F$ & $F_\eth$ & $F_\hbar$ & promotion \\
$K$ & $K_\text{mod}$ & $K_\text{slow}$ & promotion \\
$\Gamma$ & $\Gamma_\text{broad}$ & $\Gamma_\text{seq}$ & \textbf{inversion} \\
$H$ & $H_1$ & $H_\infty$ & promotion \\
$S$ & $n{:}n$ & $n{:}m$ & promotion \\
$\Omega$ & $\Omega_{\bZ_2}$ & $\Omega_\bZ$ & promotion \\
\bottomrule
\end{tabular}
\end{center}

The $\Gamma$-inversion deserves comment. The grammar encodes $\Gamma_\text{broad}$
(ordinal 4, the broadest interaction grammar), while IUG encodes $\Gamma_\text{seq}$
(ordinal 3). IUG is a restriction of the grammar on this primitive, not an extension.
This is structurally expected: a type theory that serves as the floor for many systems
must have broader interaction grammar than any single system it supports. The grammar
exceeds IUG's interaction breadth not despite being its floor, but because it is.

The lattice relations are:
$\text{meet}(\text{IUG},\,\text{grammar}) = \text{grammar}$
and $\text{join}(\text{IUG},\,\text{grammar}) = \text{IUG}$.
The grammar is lattice-contained in IUG. It is IUG's Frobenius floor. The 7-step path
above --- 6 upward promotions plus one $\Gamma$-restriction --- is the minimal additional
structure needed to reach IUG from the floor. Every step is implementable within the
grammar's existing $O_\infty$ foundation.

Note too that HoTT is not a substitute. HoTT derives bulk-boundary equivalence from the
univalence axiom; HTT requires it as a primitive. Univalence gives $\mu \circ \delta \cong \text{id}$;
it does not give $\mu \circ \delta = \text{id}$. The difference between isomorphism and
identity is precisely what the Mochizuki--Scholze--Stix dispute is about.

\subsection{A constructive sequence}
\label{subsec:constructive}

The retrosynthetic decomposition of IUG's encoding gives the required build order for any
proof assistant aiming to natively express the theory. The sequence, from most to least
structurally distinctive, is $T_\odot \to D_\odot \to R_\text{lr} \to H_\infty$.
Attempting the steps out of order --- for instance, building the $F_\ell \leftrightarrow F_\hbar$
translation layer before the $T_\odot$ foundation exists --- produces a translation layer
that cannot be grounded.

The four remaining steps for a complete verification interface are: (Step 1) identify the
$F_\hbar$ steps in IUG explicitly, construct decoherence maps showing what classical
information is lost at each, and specify what additional axiom would recover it. This is a
finite, well-defined task: the $F_\hbar$ steps are precisely the steps Scholze--Stix cannot
follow, and they are already localized. (Step 2) complete --- the grammar is the realized
HTT foundation. (Step 3) express $\Gamma_\text{seq}$ dependencies as type-theoretic temporal
dependencies; linear type theory or session types from the concurrency literature provide
the relevant technology. (Step 4) verify Corollary 3.12 within the extended system.

With Steps 1, 3, and 4 still open, the above is a program, not a proof. But the foundational
platform for executing it is not hypothetical.

\section{Predictions}
\label{sec:predictions}

Five predictions follow from the structural analysis. All are falsifiable by outcomes
independent of any encoding choice.

\paragraph{P-112 (Confirmed, Tier I).}
The Scholze--Stix impasse is a $\Phi_c + \Omega_\bZ$ structural incompatibility, not a
locatable mathematical error. The specific prediction is that the controversy will not be
resolved by either party identifying a specific, locally correctable step: a classical
$\Gamma_\text{and}$-type disagreement about an isolable line. Resolution, if it comes,
requires constructing the $F_\hbar \leftrightarrow F_\ell$ translation layer of Step 1.
\textit{Falsification condition:} the controversy resolves by identification of a specific
wrong line. \textit{Status:} Consistent with observation. No error has been located since 2018.
No smooth reformulation has been accepted by both parties.

\paragraph{P-113 (Testable, Tier I).}
A serious formalization effort targeting IUG in Lean 4, Coq, or Isabelle/HOL will achieve
partial success: sub-theorems at the $D_\triangle$ level (elliptic curve machinery, Galois
group results, number field arithmetic) will formalize within existing Mathlib. The
$\Theta$-link construction and Corollary 3.12 will not, and will require a new primitive or
axiom with no counterpart in Mathlib's set-theoretic foundations. The formalization will
stall at the $D_\odot + T_\odot$ boundary.
\textit{Falsification condition:} IUG fully formalized in a ZFC-based proof assistant with
no new foundational axioms beyond standard Lean/Coq/Isabelle primitives.

\paragraph{P-114 (Consistent, Tier II).}
IUG's $\Gamma_\text{seq}$ structure predicts that Mochizuki's documented pedagogical
practice --- insisting the papers be read in order, declining to discuss sub-results in
isolation, characterizing the verification failure as a refusal to read rather than an error
in the text --- is not stubbornness but structural honesty. A $\Gamma_\text{seq}$ proof cannot
be verified by extracting and checking its pieces; the pieces do not have independent truth
values, only sequential validity.
\textit{Falsification condition:} Corollary 3.12 is extracted and verified as a self-contained
$\Gamma_\text{and}$ argument, without the inter-universal sequential context, while still
implying abc.

\paragraph{P-115 (Structural theorem, Tier I).}
Any proof assistant that successfully verifies the $\Theta$-link identification will encode
as $O_\infty$ upon structural analysis. Any system below $O_\infty$ that appears to formalize
the identification has formalized $\mu \circ \delta \cong \text{id}$, not $\mu \circ \delta = \text{id}$,
and will lose the Diophantine content that depends on exactness. The 7-step path from the
grammar to IUG identifies precisely what additional structure must be implemented above
the $O_\infty$ floor.
\textit{Falsification condition:} a non-$O_\infty$ type theory verifies the $\Theta$-link
identification with no new foundational primitives encoding $P_\pm^\text{sym}$ directly.

\paragraph{P-116 (New, Tier I).}
No bridging paper between the IUG-native and classical-foundations communities will receive
bilateral endorsement from both Mochizuki's group and the Scholze--Stix camp. Papers
achieving classical readability will do so by reducing the $\Theta$-link to a
$D_\triangle$-level object, stripping the $D_\odot + T_\odot$ structure; these will be
accepted classically and rejected as distortions by IUG-literate readers. Papers preserving
the $D_\odot + T_\odot$ structure will be rejected classically as restatements of the
original difficulty in different notation. The intersection of \{faithful\} and
\{classically transmissible\} is empty, because the $\Omega_\bZ$ obstruction is topological:
no smooth homotopy connects $D_\odot$ to $D_\triangle$.
\textit{Falsification condition:} a single bridging paper receives explicit endorsement
from both (a) Mochizuki's group and (b) Scholze, Stix, or equivalent.
\textit{Status:} Untested in the positive direction. No bridging paper has appeared;
the 2018 Scholze--Stix report is itself informative: intended as a localized objection,
Mochizuki's response treats it as a symptom of a globally missing framework --- consistent
with $\Phi_c$ amplification and outcome (A).

\section{Open Problems}
\label{sec:open}

We list these not as a checklist but as an attempt to be honest about exactly where the
analysis is incomplete.

\begin{enumerate}

\item \textbf{Derive $P_\pm^\text{sym}$ from IUG's axioms, or show it cannot be derived.}
The $O_\infty$ assignment is a structural diagnosis, not a proved theorem about IUG. The
most direct path to confirmation is a derivation from IUG's axioms that $\mu \circ \delta = \text{id}$
exactly. The most direct path to falsification is a proof that only $\mu \circ \delta \cong \text{id}$
holds --- and a separate argument that this is sufficient for the Diophantine bound. Either
result would substantially sharpen the analysis.

\item \textbf{Build the $F_\hbar \leftrightarrow F_\ell$ translation layer.}
Step 1 of the constructive proposal is finite and well-defined, but no one has done it.
The $F_\hbar$ steps in IUG are precisely those that Scholze--Stix cannot follow; they are
already localized by the 2018 report. Each one needs a decoherence map and a specification
of what axiom would recover the lost information classically.

\item \textbf{Implement $\Gamma_\text{seq}$ dependencies in a type-theoretic setting.}
Session types and linear type theory provide the relevant technology. What is missing is the
specific implementation that expresses the $\Theta$-link step as a session-typed protocol:
``given that Copy 1 has been constructed with the following specific properties, Copy 2 is
now constructed, and the identification is now valid.'' The session type of the protocol
would be the direct answer to the Scholze--Stix objection.

\item \textbf{Check whether the $P_\pm^\text{sym}$ assignment propagates to a characterization
of which minds can generate IUG.} The analysis of \S\ref{sec:resolution} establishes a
structural necessary condition: any verification system for IUG must encode $O_\infty$
directly. The analogous question for human cognition --- what structural capacities are
necessary to generate IUG rather than merely formalize it --- is interesting and we have
not pursued it carefully enough to say much. We note only that the retrosynthetic sequence
$T_\odot \to D_\odot \to R_\text{lr} \to H_\infty$ is a structural cognitive map, not
merely a build order for a proof assistant.

\end{enumerate}

\section*{Conclusion}

The IUG controversy is a signal. Its specific sociological form --- eight years, no
locatable error, no bilateral acceptance, both parties reading the other as unreasonable ---
is not the form of a wrong proof waiting to be corrected. It is the form of a correct
transmission failure. The grammar gives this a precise statement: the impasse is the
structurally necessary consequence of $\Phi_c + \Omega_\bZ$ acting on a theory that must
be transmitted through an $F_\ell$ channel.

We have not settled whether IUG is correct. That is not what this paper is about. What we
have done is locate the obstruction precisely enough to specify what would be required to
resolve it. The obstruction is the Frobenius cliff: a non-synthesizable structural condition
that no classical accumulation process reaches. The resolution is to start on the other
side of the cliff, which requires a foundational system that already encodes $P_\pm^\text{sym}$
as a primitive. Such a system exists. What remains is the work of actually using it.

The correct response to the IUG crisis is not to ask IUG to become classical. It is to
build the interfaces that classical mathematics currently lacks. That is a mathematical
task, not a sociological one.

\vspace{1em}
\noindent\textbf{Acknowledgments.}
We thank Mochizuki for producing something structurally unusual enough to be interesting,
and Scholze and Stix for writing down the objection precisely enough to be analyzed.
This work was supported by the author's independent research.

\section*{Declarations}

\noindent\textbf{Competing Interests.}
The author has no relevant financial or non-financial interests to disclose.

\noindent\textbf{Data Availability.}
This manuscript presents a structural analysis and does not rely on or generate empirical
data. No datasets were created or analyzed.

\noindent\textbf{Funding.}
This work was supported by the author's independent research program.

\begin{thebibliography}{9}

\bibitem{synthref}
Mills, L.: SynthOmnicon (2026). \url{https://github.com/umpolungfish/synthomnicon}

\bibitem{mochizuki2012}
Mochizuki, S.: Inter-universal Teichm\"uller Theory I--IV. RIMS preprints (2012, revised 2021)

\bibitem{scholze-stix2018}
Scholze, P., Stix, J.: Why $abc$ is still a conjecture. Preprint (2018).
\url{https://www.math.uni-bonn.de/people/scholze/WhyABCisStillaConjecture.pdf}

\bibitem{mochizuki-reply2018}
Mochizuki, S.: Report on discussions, meetings, and email correspondences
(March--May 2018 with Scholze and Stix). RIMS preprint (2018)

\end{thebibliography}

\end{document}
